\clearpage\begingroup\NoSourceHead\SourceLeaf{505b}
\renewcommand{\thefootnote}{\arabic{footnote})}\setcounter{footnote}{70}
\ArithmeticHead{48}{L}{BOLZA, GAUSS UND DIE VARIATIONSRECHNUNG.}
\noindent\hspace*{\parindent}26. Wir haben in diesem Zusammenhang noch eine eigenartige und
elegante Methode zu erwähnen, die H. \textsc{Minkowski} (1907) in seinem Encyklopädieartikel\footnote{\textit{Encyklopädie der Mathematischen Wissenschaften}, Bd. V 1, S. 561—563.}
über Kapillarität zur Berechnung von $\delta U$ angewandt hat.

Der \textsc{Euler}schen Auffassung des Variationsprozesses entsprechend wird die
Fläche $U$ als Individuum einer Schar von Flächenstücken mit den Flächenparametern
$u,v$ und dem Scharparameter $\varepsilon$ aufgefasst, wobei alle Flächenstücke
demselben Bereich in der $u,v$-Ebene entsprechen. Als Parameterkurven werden
nun auf allen Flächen der Schar die beiden Scharen von Krümmungskurven
gewählt. Die Variabeln $u,v,\varepsilon$ lassen sich dann auch als krummlinige Raumkoordinaten
auffassen; für diese lässt sich das Quadrat des Linienelements
auf die Form bringen
\begin{equation*}\tag{65}
dx^2+dy^2+dz^2=L_1^2(du-l_1\,d\varepsilon)^2+L_2^2(dv-l_2\,d\varepsilon)^2+N^2\,d\varepsilon^2.
\end{equation*}

Für irgend eine Fläche der Schar lautet dann der Ausdruck für den
Flächeninhalt
\[
U(\varepsilon)=\iint L_1 L_2\,du\,dv,
\]
und daher ist
\begin{equation*}\tag{66}
\frac{dU(\varepsilon)}{d\varepsilon}=\iint\left(L_1\frac{\partial L_2}{\partial\varepsilon}+L_2\frac{\partial L_1}{\partial\varepsilon}\right)du\,dv.
\end{equation*}

Zur Umformung des Integranden wird nunmehr durch eine geometrische
Infinitesimalbetrachtung, bei der von der charakteristischen Eigenschaft der
Krümmungslinien Gebrauch gemacht wird, und in deren Mittelpunkt die
\clearpage\SourceLeaf{506}
\ArithmeticHead{49}{R}{WIRKUNGEN DER GAUSSSCHEN ARBEIT.}
\noindent \textsc{Bertrand}sche Proportionalität (63) steht, für die Hauptkrümmungsradien die
Relation abgeleitet
\[
\frac{N}{R}=\frac{1}{L_1}\left(\frac{\partial L_1}{\partial u}l_1+\frac{\partial L_2}{\partial v}l_2+\frac{\partial L_1}{\partial\varepsilon}+L_1\frac{\partial l_1}{\partial\varepsilon}\right)
\]
und eine analoge für $R'$. Aus beiden zusammen ergibt sich dann der folgende
Ausdruck für die mittlere Krümmung
\begin{equation*}\tag{67}
L_1 L_2 N\left(\frac{1}{R}+\frac{1}{R'}\right)
=L_1\frac{\partial L_2}{\partial\varepsilon}+L_2\frac{\partial L_1}{\partial\varepsilon}
+\frac{\partial(L_1 L_2 l_1)}{\partial u}+\frac{\partial(L_1 L_2 l_2)}{\partial v}.
\end{equation*}

Setzt man jetzt in (66) für den Integranden den aus dieser Gleichung
sich ergebenden Wert ein und macht von der \textsc{Green}schen Formel Gebrauch,
so erhält man nach leichter Rechnung die \textsc{Gauss}sche Formel (45) für $\delta U$.

Übrigens hatte bereits H. A. \textsc{Schwarz}\footnote{\textit{Ueber ein die Flächen kleinsten Flächeninhalts betreffendes Problem der Variationsrechnung} (1885),
Gesammelte mathematische Abhandlungen, Bd. I, S. 229—233; vgl. auch unten Nr. 27.} unter etwas anderen spezialisierenden
Voraussetzungen über die krummlinigen Koordinaten $u,v,\varepsilon$ die der
\textsc{Minkowski}schen Formel (67) entsprechende Formel rein analytisch abgeleitet.
Er setzt allgemein
\begin{equation*}\tag{68}
dx^2+dy^2+dz^2=E\,du^2+2F\,du\,dv+G\,dv^2+2H\,du\,d\varepsilon+2I\,dv\,d\varepsilon+K\,d\varepsilon^2
\end{equation*}
und macht dann die spezialisierende Annahme, dass $H=0$, $I=0$. Unter
dieser Voraussetzung beweist er die Gleichung
\begin{equation*}\tag{69}
\sqrt{EG-F^2}\sqrt K\left(\frac{1}{R}+\frac{1}{R'}\right)
=\frac{\partial\sqrt{EG-F^2}}{\partial\varepsilon}.
\end{equation*}

Sowohl die \textsc{Minkowski}sche Formel (67) als die \textsc{Schwarz}sche Formel (69)
sind übrigens als spezielle Fälle in der folgenden allgemeinen Formel enthalten,
die für beliebige Flächenscharen ohne jede spezialisierende Voraussetzung
unter Annahme des Quadrats des Linienelementes in der Form (68)
gilt:
\begin{equation*}\tag{70}
J\left(\frac{1}{R}+\frac{1}{R'}\right)
=\frac{\partial}{\partial\varepsilon}\sqrt{EG-F^2}
-\frac{\partial}{\partial u}\left(\frac{GH-FI}{\sqrt{EG-F^2}}\right)
-\frac{\partial}{\partial v}\left(\frac{EI-FH}{\sqrt{EG-F^2}}\right).
\end{equation*}
Dabei bedeutet $J$ die Funktionaldeterminante
\[
J=\frac{\partial(x,y,z)}{\partial(u,v,\varepsilon)}.
\]
\par\smallskip\noindent{\small X 2 Abh. 5.}\hfill{\small 7}\par
\clearpage\endgroup
